\section{A fixed six-term field with 32 cycles}\label{six:fixed}
The cleared and coprime representatives and their actual degrees are related in Section~\ref{sup:primitive-fields}.
In Section~\ref{se:sec:separated} take \(p=2,\ k=1,\ d=5,\ \delta=t=2^{-4096},\ K=10^6,\ \lambda=t^{12},\ \epsilon=t^{100}\). In \eqref{se:density} use \(q=1/16\) and zero constant shift. In the collected field \eqref{se:poly}, \(b_1=B=\delta\) and all four vertical rows are retained. The common density and the single regular scale below give eight simple returns per quadrant. Thus \(E_6=\max\{i+j:[x^iy^j]P\ne0\ \text{or}\ [x^iy^j]Q\ne0\}\) for the field \eqref{se:poly} at these parameters. Its binary expansion has 524,615
\unskip\ digits. The six exponent pairs determine this integer exactly. The companion archive \cite{LuReplay} contains the complete integer and the bracket data.
\subsection{The eight model brackets}
Let \(A\) be the unshifted quintic just specified. Its coefficients have
modulus below 32, and \(\sum |[w^j]A|3^j<2^{14}\).
Use the centers \(a_j=N_j/2^{380}\) in Table~\ref{six:centers} and radius
\(\sigma=2^{-330}\). Each bracket is contained in
\((2^{-81},1-2^{-97})\). In increasing order, the signs of \(J[A]\) at
the left endpoints alternate starting with plus. The right endpoints
have opposite signs. All sixteen absolute values exceed \(2^{-465}\).

To reproduce this comparison for \(h\le1/2\), write
\[
 J_j(h)=\frac{4h^{j+3/2}}{(2j+1)(2j+3)}
 {}_2F_1(1,1;j+5/2;h).
\]
This follows by expanding \((1-s)^{-1}\) in \eqref{se:mom} and using the
beta integral. Starting with one, the term ratio at index \(l\) is
\(h(2l+2)/(2l+2j+5)\). Round each term outwards on the grid
\(2^{-640}\mathbb Z\) through index 640 and add \([0,2^{-640}]\) for the
tail. Bound each square root by the enclosing dyadic endpoints on the
same grid, by integer squaring. For the upper brackets \(v=1-h<2^{-20}\).
Combine \eqref{se:explicit} into
\[
 J[A](h)=\sqrt{1-v}\,P(v)+\sqrt v\,\arccos\sqrt v\,Q(v)
\]
with rational polynomials \(P,Q\) of degree at most five. In the identity
\[
 \sqrt v\arccos\sqrt v=\tfrac{\pi}{2}\sqrt v-
       v\sum_{l\ge0}\frac{\binom{2l}{l}}{4^l(2l+1)}v^l,
\]
truncate the sum at 32 with the same grid and tail bound.
Machin's formula \(\pi/2=8\arctan(1/5)-2\arctan(1/239)\), using
160 alternating terms and the first omitted term, encloses the remaining
constant. These operations give the stated strict inequalities.

\begin{table}[ht]
\centering\scriptsize
\begin{tabular}{@{}ll@{}}\toprule
\multicolumn{2}{c}{Hexadecimal numerator \(N_j\), most significant part on the left}\\ \midrule
\texttt{0000000000000004fffffdeeb8a56c7560e9f748c8991d4} & \texttt{5306200a13f9a6fd064824218d0fbba805e052f39be7f77a}\\
\texttt{0000000254f631822ceb3916f44886e1960d86c9d679a4f} & \texttt{957b15c83f0d02696092bbfe59554ecbd75a9a6654efa73d}\\
\texttt{678aef2419e14dddd10f9792d94af7090d420c65a604f2d} & \texttt{c9f3adbb20fa0afe864507621d80c0fb6e3d93c89e54cfd1}\\
\texttt{ffffff6132ab842bb50af286544f9158d20b3d8ee72cea0} & \texttt{9ba85bc64a44f5101043dab8f4aba93f05ee54827fff721b}\\
\texttt{fffffffffb876590a5021c23c445ad8e3cd7712f7e12a7a} & \texttt{9f844ada5ff696864c277a009edaf72a1a9e76f6aee915b7}\\
\texttt{ffffffffffffffcbb7cc09bc6cfea7fde8c3d6464a90aea} & \texttt{6036944009383c610585ddcced735556346669fb09670156}\\
\texttt{ffffffffffffffffd3dc857a7a9d51689139de77054cb72} & \texttt{585b8eadcf395769ba6215cb3a2f23cc8c6ee563432772d7}\\
\texttt{fffffffffffffffffffffff983a008ec5c452f18e381d54} & \texttt{9125a7fe770685caa7c746f7eaa55e237cebbd4e8f70fc30}\\
\bottomrule
\end{tabular}
\caption{Concatenate each pair of strings and divide the resulting integer by
\(2^{380}\). All eight brackets have rational endpoints.}\label{six:centers}
\end{table}

On the union of the brackets \(|J[A]''|<2^{180}\).
For \(h\le1/2\), the center series without its prefactor has coefficients
at most one, so its derivatives of orders \(l=0,1,2\) are bounded by
\(l!2^{l+1}\). Differentiate the displayed prefactors using \(h>2^{-81}\)
and the coefficient bound 32. This gives less than \(2^{65}\).
For \(h\ge1/2\) the coefficient absolute sums of \(P,Q\) and their first
two derivatives are below \(2^{20}\), by rational expansion.
On \(2^{-97}<v\le1/2\), \(g(v)=\sqrt v\arccos\sqrt v\) has
\(|g'|\le 2/\sqrt v,\ |g''|\le4/v^{3/2}\) by differentiation.
These estimates give the stated bound. At some point of each bracket
the secant slope has magnitude greater than \(2^{-465}/\sigma\).
Its variation in the bracket is less than \(2\sigma\,2^{180}\), hence
\begin{equation}
 |J[A]'|>2^{-136}\quad\hbox{throughout each bracket}.              \label{six:margins}
\end{equation}

\subsection{Finite horizontal exponents}
Use the single index \(i=1\), \(E=h\), \(b_i=B=\delta\), \(W_i(x)=W_\delta(-x/b_i)=\delta(e^{-x/\delta}-1)+1-e^{-x}\),
\(W_*(x)=i-1+1-e^{-x}\) for \(x\ge0\).
On \(\sqrt\delta\le x\le128\) we have
\[
 |W_i-W_*|<390\delta=:\eta,\qquad W_i',W_*'>2^{-193}.
\]
Indeed only the wall term differs from the model, by at most \(B=\delta\) on this interval. In the derivative only the wall term is negative.
Its magnitude is at most
\((B/b_i)e^{-\sqrt\delta/b_i}\le2e^{-1/\sqrt\delta}<2\delta\),
whereas the \(i\)-th positive term is \(e^{-x}\ge e^{-128}>2^{-192}\).
At \(x=128\) the potential exceeds \(E\), since \(e^{-128}+\eta<2^{-97}\).
At \(x=-\sqrt\delta\) it exceeds 3 by
\(W_\delta(u)\ge e^{\delta u}(\delta e^{(1-\delta)u}-3)-B\) for
\(u=\sqrt\delta/b_i\ge1/\sqrt\delta\).
Thus all the actual integrals are supported in
\(-\sqrt\delta\le x\le128,\ |Z|\le\sqrt3\).

Set \(\Phi(s)=2\int_0^{\sqrt{s_+}} A(Z^2)\dd Z\).
It is Hölder continuous there with modulus \(2^{15}\sqrt{|\Delta s|}\).
The interval \(|x|\le\sqrt\delta\) and the value error \(\eta\) on the
rest give
\[
 |b_i{\cal I}_{\delta,A}(E)-J[A](h)|<2^{30}\sqrt\delta<2^{-2000}.
\]
For its energy derivative the slice is \(A(s)/\sqrt s\) for \(s>0\)
and zero for \(s<0\), with \(s=E-W\). The turning points are simple, so
the inverse-square-root singularities are integrable.
On \(0\le x\le\sqrt\delta\) the gap exceeds \(h/2\) in both models.
On \(-\sqrt\delta\le x<0\), where \(W_i<E/2\), use \(E-W_i>E/2\).
On the remaining part of this interval with \(W_i\ge E/2\), convexity
of \(W_\delta\) for \(u\ge0\) gives \(|W_i'|\ge W_i/|x|\ge
E/(2\sqrt\delta)\). Substitution by \(W_i\) shows that the sum of all
these boundary contributions is less than \(2^{105}\sqrt\delta\).

On the main interval put \(\tau=2^{-1024}>\eta\).
Where \(|E-W_*|\le2\tau\), substitution by each increasing potential
bounds the sum of the two singular integrals by
\(2^{212}\sqrt\tau\), using \(0<E-W\le3\tau\) on their contributing
parts and the derivative bound \(2^{-193}\). Where \(E-W_*>2\tau\),
both arguments exceed \(\tau\) and
\(|(A(s)/\sqrt s)'|\le2^{20}\tau^{-3/2}\) up to \(s=3\).
The rest contributes zero. Thus on each bracket
\begin{equation}
 |b_i{\cal I}'_{\delta,A}(E)-J[A]'(h)|
 <2^{105}\sqrt\delta+2^{212}\sqrt\tau+128\cdot2^{20}\eta\tau^{-3/2}
 <2^{-298}.                                                        \label{six:deltaestimate}
\end{equation}
The same substitutions bound the full integral of
\(\mathbf1_{W_i<E}/\sqrt{E-W_i}\) by \(2^{210}\).

\subsection{Finite vertical scale}
For \(p=2\), the matrix~\eqref{se:vertical} is
\[
\begin{pmatrix}
 1&1&0&0&0&0\\
 0&1&3/2&5/2&-3&2\\
 0&0&-1/12&1/4&5/3&-7/6\\
 0&0&-1/60&1/60&0&-1/5\\
 0&0&-1/168&1/280&0&-11/140\\
 0&0&-1/360&1/840&0&-13/315
\end{pmatrix}.
\]
Its determinant is \(-1/4762800\) and its inverse maximum row sum is
\(91804/5\).
All six controls in \eqref{se:controls} have magnitude below \(6\cdot10^5\),
and the base weights are \(2,-3,1\). Here \(\lambda=t^{12}\).
On \(|Z|\le2\) the error in \eqref{se:transport} is less than
\(2^{110}t^2\). To verify it, temporarily replace \(\lambda\) by a
complex variable \(\zeta\) of modulus at most \(r_0=10^{-12}\), with
the controls fixed. On \(|z|\le1/2\), \(|\log L|<1\),
\(|L^{q_j}|<4,\ |L^\rho|<2,\ \sum_{j\ge1}|C_j|<5\cdot10^{-6}\).
Bounding the derivatives of the powers on this disk gives
\(|(L^{q_j}-L^\rho)/z|<50\). Hence \(|V'/z-1|<1/32\), and integrating
and taking the square root gives \(|Y/z-1|<1/32\).
Consequently \(|Y_z-1|<1/10\) and comparison with \(z-Y\) on
\(|z|=1/2\) gives an analytic inverse for \(|Y|\le1/4\).
There the even transported density has modulus less than 100 since
\(|z'|<2\) and \((|\gamma_0|,|\gamma_1|,|\gamma_2|)\le(3,4,1)\).
Cauchy's bound on the nonlinear remainder is \(200(\lambda/r_0)^2\),
and the Taylor tail of its coefficient of \(\zeta\) at \(Y=tZ\) is at
most \(200(8t)^{12}/r_0\). Division in \eqref{se:transport} proves the
claim. Integrating the error and its boundary values, using the slice
derivative formula and the last bound in the preceding subsection,
gives value and derivative errors at most \(2^{320}t^2<2^{-2000}\)
for \(b_i M/(\lambda t^{11})\) relative to \(b_i{\cal I}_{\delta,A}\).

\subsection{The full return at the specified perturbation}
The following uniform complex neighborhood estimates control the
higher orders of the return. They apply at every real \(E_0=i-1+h\)
in the eight brackets. In the horizontal coordinate
\(w=\operatorname{sgn}(u)\sqrt{W_\delta(u)}\) write \(u=u(w)\).
For real \(|w|\le\sqrt{E_0}\) we have
\(-128/b_1<u<8192\). The left bound follows from the single exponential estimate at 128, and the right bound from
\(W_\delta(8192)>3\). In particular \(|u|<2^{12300}\).
We show that around each such real \(w_0\) its inverse extends to
\begin{equation}
          |w-w_0|<2^{-97000},\qquad |u(w)|<2^{12301}.               \label{six:inverse}
\end{equation}

Write \(W_\delta(u)=u^2 H_1(u)\), where
\(H_1(u)=\int_0^1(1-s)W_\delta''(su)\dd s\). Along the real segment
from 0 to \(u_0=u(w_0)\) we have \(B e^u<5\) by \eqref{se:W},
\(W_\delta(u)\le3\) and \(u<8192\). On its unit complex neighborhood
\(|W_\delta''|,|W_\delta'''|<32\). Hence \(|H_1|,|H_1'|<32\)
on the unit disk around \(u_0\). For real \(|u|\le1\),
\(W_\delta''>B/4\). For \(|u_0|\ge1\),
\(W_\delta(u_0)\ge B/8\) by monotonicity. It follows in both cases that
\(H_1(u_0)>2^{-28700}\). On \(|u-u_0|\le2^{-29000}\) the real-positive
square-root branch continues and \(|u\sqrt{H_1(u)}|<2^{12305}\).
On the disk of half that radius its second derivative has modulus
less than \(2^{70400}\) by Cauchy's estimate.

The real \(w'(u_0)\) exceeds \(2^{-12500}\). For \(|u_0|\le1\) use
\(W_\delta''>B/4\) and \(w'=W_\delta'/(2u\sqrt{H_1})\) with the
continuous value at zero. For \(u_0\ge1\) use convexity and
\(W_\delta'(1)>B/4\). For \(u_0\le-1\), one negative term alone gives
\[
 |W_\delta'(u_0)|\ge b_1 e^{-128}(1-e^{-(1-b_1)})>2^{-12482}.
\]
Divide by \(2\sqrt{W_\delta(u_0)}<4\). On \(|u-u_0|\le2^{-84000}\)
the variation of \(w'\) is less than a quarter of \(w'(u_0)\).
Comparison with its linear part on the boundary gives a unique analytic
inverse on the disk in \eqref{six:inverse}. The local inverses agree along
the real interval.

After an increasing change of time, the \((u,Z)\) field is
\(\dot u=-2Z+A_\varepsilon,\ \dot Z=W_\delta'(u)\), with
\[
 A_\varepsilon=\frac{2}{t L Y_z}
          \big(\dot u_{\rm original}-(\dot u_{\rm original})|_{\varepsilon=0}\big).
\]
Thus in \(w=\sqrt E\cos\theta,\ Z=\sqrt E\sin\theta\) the angular
equation for a turn is exactly
\begin{equation}
              \frac{dE}{d\theta}=\frac{2wA_\varepsilon}{2-ZA_\varepsilon/E}.
                                                                        \label{six:angle}
\end{equation}
The factor \(w'(u)>0\) cancels. Put \(R=2^{-100000}\) and
\(r_\varepsilon=2^{-117000}\).
For \(|E-E_0|\le4R\) and real \(\theta\), the principal square root
varies by at most \(2^{42}|E-E_0|\), by \(E_0>2^{-81}\).
Hence \eqref{six:inverse} applies and \(|Z|,|w|<2\).
For \(|\varepsilon|\le r_\varepsilon\) all exponential arguments
\(\varepsilon\gamma_j u/C_j\) have modulus less than \(1/2\), since
\(|C_0|>1/2,\ |C_{1,2}|>\lambda\).
The vertical estimates give \(|1/(LY_z)|<4\) and both types of vertical
powers in the perturbation have modulus below 8. Applying
\(|e^\xi-1|\le2|\xi|\) termwise gives
\(|A_\varepsilon|<2^{16415}|\varepsilon|\) for nonzero \(\varepsilon\).
The denominator in \eqref{six:angle} has modulus greater than one and
\(|dE/d\theta|\le2^{16420}|\varepsilon|\). A solution starting within
\(2R\) of \(E_0\) changes by less than \(R\) during a full turn, since
\(8\cdot2^{16420} r_\varepsilon<R\).
Its displacement \(\Delta\) is therefore holomorphic in the initial
\(E\) and in \(\varepsilon\) there, with \(|\Delta|<R\).

The linear coefficient of \(\Delta\) is \(2M/t^2\). For the stated
\(\varepsilon=t^{100}<r_\varepsilon/2\), Cauchy's estimates in the two
variables bound the value and \(E\)-derivative of
\(\Delta-2\varepsilon M/t^2\) at \(E_0\) by
\(2(\varepsilon/r_\varepsilon)^2\). In the normalized displacement
\(b_i\Delta/(2\varepsilon\lambda t^9)\) the resulting errors are less than
\[
            \frac{2\varepsilon}{t^{21}r_\varepsilon^2}<2^{-1000}.
\]
Together with \eqref{six:margins}, \eqref{six:deltaestimate} and the value
bounds, this preserves opposite endpoint signs and a strict derivative
sign on all eight real intervals. The angular speed is positive there.
Every interval therefore gives a simple fixed point of the actual
one-turn map, with positive multiplier different from one.
The compact, disjoint radial brackets on the section distinguish
the cycles, and uniqueness for the scalar angular equation orders
the corresponding closed graphs. Here \(\rho,\ q_j-N_j,\alpha_j\)
all have modulus below \(1/4\) for the specified parameters.
The support check in \eqref{se:poly} applies exactly. The four quadrants
give 32 hyperbolic cycles at the stated parameters.
